\section{First law of thermodynamics}

	The first law is the statement of conservation of energy applied to thermodynamic systems: energy can be transferred as heat or work and internal energy changes accordingly. For a closed system (fixed mass)
	\begin{align*}
		\Delta U&=Q-W
	\end{align*}
	where $Q$ is heat supplied to the system and $W$ is work done by the system. In differential form
	\begin{align*}
		\dd{U}&=\delta{Q}-\delta{W}
	\end{align*}
	\begin{itemize}
		\item For reversible PV-work, we have $\delta{W}=P\dd{V}$. Hence
		\begin{align*}
			\delta{Q}&=\dd{U}+P\dd{V}
		\end{align*}
		
		\item $Q$ and $W$ are path-dependent (inexact differentials, often written as $\delta{Q}$, $\delta{W}$).
		
		\item $U$ is a state function (exact differential, written as $\dd{U}$).
		
		\item First law does not constrain the direction of processes.
	\end{itemize}
	
	\subsection{Applications of the first law of thermodynamics}
	
		\subsubsection{General relation between $C_P$ and $C_V$}
			Specific heats are defined for closed systems of fixed composition:
			\begin{align*}
				C_V\equiv\qty(\pdv{U}{T})_V\qq{and}C_P\equiv\qty(\pdv{H}{T})_P
			\end{align*}
			where enthalpy $H\equiv U+PV$. For a general substance the difference $C_P-C_V$ can be written in terms of measurable derivatives. One convenient form is:
			\begin{align*}
				C_P-C_V&=-T\dfrac{\qty(\pdv*{P}{T})_V^2}{\qty(\pdv*{P}{V})_T}
			\end{align*}
			This identity is derived from total differentials and Maxwell relations; it is valid for any simple compressible substance.
			
			\begin{itemize}
				\item For an ideal gas, we have
				\begin{align*}
					P&=\dfrac{RT}{V_m} \\
					\implies \qty(\pdv{P}{T})_V=\dfrac{R}{V_m}&\qq{and}\qty(\pdv{P}{V})_T=-\dfrac{RT}{V_m^2}
				\end{align*}
				
				Substituting the above gives us
				\begin{align*}
					C_P-C_V&=-T\dfrac{R^2/V_m^2}{\qty(-RT/V_m^2)} \\
					\implies\Aboxed{C_P-C_V&=R}
				\end{align*}
				which is the \textit{Mayer's relation}.
			\end{itemize}